% !PRINTED 586
\weberpage{586}
Here $F_\nu$, $\psi_\nu$ are integral functionals of the absolute domain of rationality,
while $f_\nu$ are integral functionals in $R$. The expressions for $F_\nu, f_\nu, \psi_\nu$
result from the formulae \S.\,158, (15). To these functionals one can apply the same
inferences that led to the Theorem \S.\,158, 4., and it results that:

\weq{12}{
\begin{aligned}
\mathfrak{S}_R \frac{f_\nu(\tau)}{f'(\tau)} &= 0, & \nu &= 0, 1, \ldots, n-2.\\
\mathfrak{S}_R \frac{f_{n-1}(\tau)}{f'(\tau)} &= 1.
\end{aligned}
}
If now one divides the functions $F_\nu(t)$, when they are of degree higher than $(n-1)$,
by $f(t)$ and takes the remainder of the division; and then in this remainder expresses
the powers of $t$ [according to \S.\,158, (15), applied to $f_\nu$] by the functions
$f_\nu(t)$, there results the representation

\weq{13}{
F_\nu(\tau) = \alpha_0 f_0 + \alpha_1 f_1 + \cdots + \alpha_{n-1} f_{n-1},
}
in which the $\alpha_0, \alpha_1, \ldots, \alpha_{n-1}$ are \gesperrt{integral functionals} in $R$. This
representation holds for every root $\tau$ of $f(t) = 0$, that is, for $\tau = \tau_1, \tau_2,
\ldots, \tau_n$. From this, however, there results by (12):
\[
\mathfrak{S}_R \left( \frac{F_\nu(\tau)}{f'(\tau)} \right) = \alpha_{n-1}.
\]
If one here divides by $\psi'(\sigma)$ and forms the trace $S_R$, one obtains from (7),
with application of Theorem 4., \S.\,158:
\weq{14}{
S_\Omega \left( \frac{F_\nu(\tau)}{f'(\tau)\,\psi'(\sigma)} \right) = a_\nu,
}
in which $a_\nu$ denotes an integral rational functional.

By (9) we now have, for every root $\tau_i$ of $F(t) = 0$ with the exception of $\tau_i = \tau$,
\[
\tau^{n\,m-1} F_0(\tau_i) + \tau^{n\,m-2} F_1(\tau_i) + \cdots + F_{n\,m-1}(\tau_i) = 0,
\]
and
\[
\tau^{n\,m-1} F_0(\tau) + \tau^{n\,m-2} F_1(\tau) + \cdots + F_{n\,m-1}(\tau) = F'(\tau),
\]
so that, when one multiplies (14) by $\tau^{n\,m-\nu-1}$ and takes the sum with respect
to $\nu$,
\weq{15}{
\frac{F'(\tau)}{f'(\tau)\,\psi'(\sigma)} = a_0 \tau^{n\,m-1} + a_1 \tau^{n\,m-2} + \cdots
}
results. \gesperrt{Therefore $F'(\tau)$ is divisible by $f'(\tau)\,\psi'(\sigma)$.}
